\begin{table*}[t]
\centering
\caption{XGBoost router feature importance (Phase~5, 25 features). Gain = total loss reduction from splits on each feature, normalized to sum to 1. Mean $\pm$ std across 4 species (200 trees each, GroupShuffleSplit 80/20 by protein).}
\label{tab:s3_feature_importance}
\scriptsize
\setlength{\tabcolsep}{2pt}
\begin{tabular}{rllrrrrr}
\toprule
Rank & Feature & Category & Mean Gain & Mouse & Dog & Chicken & Zebrafish \\
\midrule
1 & G: on primary chrom & Geometry & 0.2113 & 0.2220 & 0.2012 & 0.2616 & 0.1606 \\
2 & A: miniprot hit & A-stage & 0.1186 & 0.0575 & 0.2139 & 0.1945 & 0.0083 \\
3 & G: strand match & Geometry & 0.0756 & 0.0288 & 0.0556 & 0.0637 & 0.1541 \\
4 & A: \#tools hitting & A-stage & 0.0614 & 0.0964 & 0.0638 & 0.0146 & 0.0707 \\
5 & A: diamond hit & A-stage & 0.0448 & 0.0198 & 0.0403 & 0.0742 & 0.0448 \\
6 & G: co-orth overlap (bp) & Geometry & 0.0434 & 0.0468 & 0.0481 & 0.0422 & 0.0366 \\
7 & C: composite score & C-stage & 0.0417 & 0.0222 & 0.0330 & 0.0275 & 0.0840 \\
8 & A: \#raw hits & A-stage & 0.0402 & 0.0883 & 0.0160 & 0.0269 & 0.0297 \\
9 & C: hit score & C-stage & 0.0353 & 0.0330 & 0.0247 & 0.0411 & 0.0422 \\
10 & S: query type (human/sp) & Species ind. & 0.0348 & 0.0412 & 0.0233 & 0.0285 & 0.0460 \\
11 & C: \#exons & C-stage & 0.0347 & 0.0477 & 0.0303 & 0.0273 & 0.0334 \\
12 & C: CDS length & C-stage & 0.0317 & 0.0315 & 0.0211 & 0.0288 & 0.0453 \\
13 & C: coverage & C-stage & 0.0300 & 0.0305 & 0.0234 & 0.0377 & 0.0286 \\
14 & C: identity & C-stage & 0.0266 & 0.0281 & 0.0414 & 0.0171 & 0.0199 \\
15 & Q: query length (aa) & Query & 0.0227 & 0.0383 & 0.0173 & 0.0153 & 0.0197 \\
16 & G: co-orth same chrom & Geometry & 0.0217 & 0.0178 & 0.0205 & 0.0235 & 0.0250 \\
17 & G: region size (bp) & Geometry & 0.0208 & 0.0191 & 0.0251 & 0.0135 & 0.0254 \\
18 & G: region size (log) & Geometry & 0.0194 & 0.0167 & 0.0242 & 0.0144 & 0.0222 \\
19 & A: best qcov & A-stage & 0.0188 & 0.0126 & 0.0132 & 0.0119 & 0.0373 \\
20 & Q: query length (log) & Query & 0.0179 & 0.0345 & 0.0175 & 0.0074 & 0.0121 \\
21 & A: best score & A-stage & 0.0172 & 0.0196 & 0.0151 & 0.0136 & 0.0203 \\
22 & A: mmseqs hit & A-stage & 0.0139 & 0.0241 & 0.0150 & 0.0077 & 0.0087 \\
23 & A: last hit & A-stage & 0.0109 & 0.0131 & 0.0049 & 0.0047 & 0.0209 \\
24 & C: has hit & C-stage & 0.0069 & 0.0101 & 0.0109 & 0.0025 & 0.0042 \\
25 & C: rank & C-stage & 0.0000 & 0.0000 & 0.0000 & 0.0000 & 0.0000 \\
\bottomrule
\end{tabular}

\vspace{3pt}
\footnotesize
\textbf{Feature categories:} A = A-stage (tool evidence), G = Geometry (genomic context), Q = Query (protein length), C = C-stage (alignment quality), S = Species indicator (0=human, 1=species protein). The top predictor is \texttt{on\_primary\_chrom} (mean gain 0.21), reflecting the router's learned spatial prior that correct paralog loci tend to reside on the same chromosome as the primary ortholog. The species indicator (\texttt{query\_type\_num}) ranks 10th (mean gain 0.035), providing moderate but consistent signal across all 4 species. The Phase~4 (human-only) feature ranking differs: \texttt{c\_identity} was the top predictor (gain 0.19) when species-protein rows were absent; in Phase~5, geometric features dominate because the dual-query setup makes alignment-quality differences more subtle within each query type.
\end{table*}
